% Preâmbulo do guião (pacotes, estilos e caixas)
% Para o relatório: copiar estas linhas para o preâmbulo do relatório.

\usepackage[utf8]{inputenc}
\usepackage[T1]{fontenc}
\usepackage[portuguese]{babel}
\usepackage{lmodern}
\usepackage{microtype}
\usepackage[a4paper,top=2.6cm,bottom=2.6cm,left=2.4cm,right=2.4cm,headheight=14pt]{geometry}
\usepackage[table]{xcolor}
\usepackage{listings}
\usepackage{tabularx}
\usepackage{array}
\usepackage{booktabs}
\usepackage{enumitem}
\usepackage{needspace}
\usepackage{graphicx}
\usepackage{tikz}
\usetikzlibrary{arrows.meta,shapes.geometric}
\usepackage{adjustbox}
\usepackage{float}
\usepackage{caption}
\usepackage{titlesec}
\usepackage{tocloft}
\usepackage{fancyhdr}
\usepackage[most]{tcolorbox}
\usepackage{hyperref}

\DeclareUnicodeCharacter{2192}{\ensuremath{\rightarrow}}
\DeclareUnicodeCharacter{2194}{\ensuremath{\leftrightarrow}}

% ---------- Cores ----------
\definecolor{azul}{RGB}{31,73,125}
\definecolor{cinza}{RGB}{90,90,90}

\hypersetup{colorlinks=true,linkcolor=azul,urlcolor=azul,pdftitle={Guião-Workshop SR2}}

% ---------- Texto ----------
\setlength{\parindent}{0pt}
\setlength{\parskip}{0.45em}
\setlist{itemsep=0.15em,topsep=0.25em}
\widowpenalty=10000
\clubpenalty=10000
\emergencystretch=3em

% ---------- Títulos e numeração ----------
% Introdução (sem número); Capítulo 1 = Parte A (1.1, 1.2...); Capítulo 2 = Parte B (2.1, 2.1.1...)
\setcounter{secnumdepth}{3}
\titleformat{\chapter}[display]{\huge\bfseries\color{azul}}{\Large\chaptertitlename\ \thechapter}{0.4em}{}[{\vspace{0.3em}\color{azul!40}\titlerule[1pt]}]
\titlespacing*{\chapter}{0pt}{0pt}{1.6em}
\titleformat{\section}{\Large\bfseries\color{azul}}{\thesection}{0.8em}{}[{\color{azul!40}\titlerule[0.8pt]}]
\titleformat{\subsection}{\large\bfseries\color{azul}}{\thesubsection}{0.7em}{}
\titleformat{\subsubsection}{\normalsize\bfseries\color{azul!80!black}}{\thesubsubsection}{0.6em}{}
\titlespacing*{\section}{0pt}{1.6em}{1.0em}
\titlespacing*{\subsection}{0pt}{1.5em}{0.6em}
\titlespacing*{\subsubsection}{0pt}{1.2em}{0.4em}

% ---------- Esquemas (desenhados em TikZ, ficheiros em esquemas/*.tex) ----------
\providecommand{\esquemadir}{esquemas/}
\counterwithout{figure}{chapter}
\captionsetup{font=small,labelfont={bf,color=azul}}
\newcommand{\ph}[1]{\texttt{\textless{}#1\textgreater{}}}
\tikzset{
  vm/.style={draw=azul, very thick, fill=azul!8, rounded corners=3pt, align=center, font=\small, inner sep=5pt},
  svc/.style={draw=azul!70, fill=white, rounded corners=2pt, align=center, font=\footnotesize, inner sep=4pt},
  grupo/.style={draw=azul!70, fill=azul!12, rounded corners=2pt, align=center, font=\footnotesize, inner sep=4pt},
  net/.style={draw=black!50, dashed, fill=black!4, rounded corners=2pt, align=center, font=\footnotesize, inner sep=4pt},
  ext/.style={draw=orange!70!black, fill=orange!10, ellipse, align=center, font=\footnotesize, inner sep=3pt},
  caixa/.style={draw=azul, thick, rounded corners=4pt, fill=azul!3},
  caixatit/.style={font=\footnotesize\bfseries, text=azul, anchor=north west},
  seta/.style={-{Stealth[length=2.2mm]}, thick, draw=black!70},
  setad/.style={seta, dashed},
  duplo/.style={{Stealth[length=2.2mm]}-{Stealth[length=2.2mm]}, thick, draw=black!70},
  linha/.style={thick, draw=black!55},
  rot/.style={font=\scriptsize, fill=white, inner sep=1.5pt, align=center},
  vida/.style={draw=black!35, dashed},
  acao/.style={draw=azul!60, fill=azul!5, rounded corners=2pt, align=center, font=\scriptsize, inner sep=3pt},
  nota/.style={draw=orange!60!black, fill=yellow!12, align=center, font=\scriptsize, inner sep=3pt},
}
\newcommand{\esquema}[2]{%
  \begin{figure}[H]
    \centering
    \adjustbox{max width=\linewidth}{\input{\esquemadir#1}}
    \caption{#2}
  \end{figure}}

% Passo de um procedimento
\newcommand{\passo}[1]{\par\needspace{6\baselineskip}\vspace{0.5em}{\color{azul}\bfseries #1}\par\nopagebreak}

% Remissão para os erros comuns
\newcommand{\problemas}[1]{\par\vspace{0.3em}{\small\color{cinza}\textit{Problemas? Ver} #1.}\par}

% Marca de secção por preencher
\newcommand{\aguardar}{{\normalfont\small\color{cinza}[A AGUARDAR CONTEÚDO]}}

% ---------- Índice ----------
\setcounter{tocdepth}{2}
\renewcommand{\cftchappresnum}{\chaptername\ }
\setlength{\cftchapnumwidth}{6.2em}
\renewcommand{\cftchapfont}{\bfseries\color{azul}}
\renewcommand{\cftchappagefont}{\bfseries\color{azul}}
\renewcommand{\cftsecleader}{\cftdotfill{\cftdotsep}}
\setlength{\cftsecindent}{1.5em}
\setlength{\cftsecnumwidth}{2.6em}
\setlength{\cftsubsecindent}{4.1em}
\setlength{\cftsubsecnumwidth}{3.4em}
\addto\captionsportuguese{\renewcommand{\contentsname}{Índice}}

% ---------- Cabeçalho e rodapé ----------
\pagestyle{fancy}
\fancyhf{}
\renewcommand{\chaptermark}[1]{\markboth{\chaptername\ \thechapter\ - #1}{}}
\renewcommand{\sectionmark}[1]{}
\fancyhead[L]{\small\color{cinza}Guião-Workshop SR2}
\fancyhead[R]{\small\color{cinza}\nouppercase{\leftmark}}
\fancyfoot[C]{\small\color{cinza}\thepage}
\renewcommand{\headrulewidth}{0.4pt}
\renewcommand{\footrulewidth}{0pt}
\fancypagestyle{plain}{\fancyhf{}\fancyfoot[C]{\small\color{cinza}\thepage}\renewcommand{\headrulewidth}{0pt}}

% ---------- Blocos de código ----------
\lstdefinestyle{base}{
  basicstyle=\ttfamily\small,
  breaklines=true,
  breakatwhitespace=false,
  columns=fullflexible,
  keepspaces=true,
  showstringspaces=false,
  frame=single,
  framesep=4pt,
  rulecolor=\color{black!20},
  backgroundcolor=\color{black!3},
  xleftmargin=4pt, xrightmargin=4pt,
  aboveskip=0.7em, belowskip=0.7em,
  upquote=true,
  literate=
    {á}{{\'a}}1 {à}{{\`a}}1 {â}{{\^a}}1 {ã}{{\~a}}1
    {é}{{\'e}}1 {ê}{{\^e}}1 {í}{{\'i}}1 {ó}{{\'o}}1
    {ô}{{\^o}}1 {õ}{{\~o}}1 {ú}{{\'u}}1 {ç}{{\c c}}1
    {Á}{{\'A}}1 {É}{{\'E}}1 {Í}{{\'I}}1 {Ó}{{\'O}}1 {Ú}{{\'U}}1 {Ç}{{\c C}}1
    {→}{{$\rightarrow$}}1
}
\lstdefinestyle{bashstyle}{style=base,language=bash,commentstyle=\color{green!40!black}\itshape}
\lstdefinestyle{psstyle}{style=base,commentstyle=\color{green!40!black}\itshape,morecomment=[l]{\#}}
\lstdefinestyle{cfgstyle}{style=base}
\lstdefinestyle{diagstyle}{style=base,frame=none,backgroundcolor=,basicstyle=\ttfamily\footnotesize}

% ---------- Caixas ----------
\tcbset{enhanced,breakable,fonttitle=\bfseries,boxrule=0.6pt,arc=2pt,left=6pt,right=6pt,top=4pt,bottom=4pt,before skip=0.8em,after skip=0.8em}
\newtcolorbox{avisobox}[1][]{colback=red!3,colframe=red!55!black,#1}
\newtcolorbox{notabox}[1][]{colback=azul!4,colframe=azul,#1}
\newtcolorbox{importantebox}[1][]{colback=orange!5,colframe=orange!70!black,#1}
\newtcolorbox{dicabox}[1][]{colback=green!4,colframe=green!40!black,#1}
\newtcolorbox{simplesbox}{colback=black!3,frame hidden,borderline west={2pt}{0pt}{black!35},sharp corners}
